\documentclass[11pt]{article}
\usepackage[margin=1in]{geometry}
\usepackage{amsmath,amssymb,amsthm,booktabs}
\usepackage[hidelinks]{hyperref}
\usepackage{xurl}
\setlength{\emergencystretch}{3em}
\newtheorem{theorem}{Theorem}
\newtheorem{corollary}[theorem]{Corollary}
\newcommand{\C}{\mathbb C}
\newcommand{\Q}{\mathbb Q}
\newcommand{\mr}{\operatorname{mr}}
\title{The symmetric-group quotient of the reflection class of $G_{31}$}
\author{Exact finite certificate and method comparison}
\date{17 November 2026}
\begin{document}
\maketitle
The sixty involutory reflections of $G_{31}$ map onto the fifteen
transpositions of $S_6$, four to one. The map is induced by the exact
matrix action on fifteen blocks of root lines, so it is a group
homomorphism. A meridional epimorphism onto this reflection class
therefore already supplies both a rank-five Coxeter quotient and a
five-dimensional irreducible KCH representation. The class bound five
does not provide stronger detection than those two existing bounds.
No new named-knot equality is asserted here.

\section*{The matrix model and quotient}
Put $K=\Q(i)$ and use the standard positive Hermitian form on $K^4$.
For a nonzero column $u$, put
\[
 s_u=I-\frac{2u\overline u^{\,T}}{\overline u^{\,T}u}.
\]
The sixty root lines are represented, in the following order, by
\[
 e_a\ (1\le a\le4),\qquad
 e_a-i^k e_b\ (1\le a<b\le4,\ 0\le k<4),
\]
followed by
\[
 (1,i^a,i^b,i^c)^T
 \quad(0\le a,b,c<4,\quad a+b+c\equiv0\pmod2).
\]
Let $C$ be their sixty reflections, and let $W$ be generated by the
reflections in the five columns
\begin{equation}\label{eq:generators}
 (0,0,0,1)^T,\quad(1,-1,0,0)^T,\quad(1,0,0,1)^T,\quad
 (1,i,1,-i)^T,\quad(1,1,1,1)^T .
\end{equation}
This is the $G_{31}$ model in Konishi--Minabe
\cite[Section~6.1, equation~(6.1)]{km25}; the retained accepted manuscript
places that equation on printed page~25. The displayed root-pair index
$0\le\alpha<\beta\le4$ is read $1\le\alpha<\beta\le4$ in its announced
basis $e_1,\ldots,e_4$. Its thirty-two four-coordinate roots are
projective lines. More precisely, the standard monomial
$G(4,2,4)$ orbit of $(1,1,1,1)^T$ consists of the 128 vectors
$(i^a,i^b,i^c,i^d)^T$ with $0\le a,b,c,d<4$ and $a+b+c+d$ even.
Dividing by the four scalar phases gives these thirty-two lines.
This count concerns the monomial subgroup, not the full $G_{31}$ orbit.

\begin{theorem}\label{thm:quotient}
The exact matrix group $W$ has order $46080$. Its reflections are
precisely $C$, and they form one conjugacy class. There is a
surjective homomorphism
\[
                    \pi:W\longrightarrow S_6
\]
with kernel of order $64$. Its restriction to $C$ is four to one onto
the transpositions. The images of the five generators in
\eqref{eq:generators}, in order, are
\[
                    (12),\ (34),\ (45),\ (14),\ (16).
\]
\end{theorem}
\begin{proof}
Number the sixty lines from $0$ to $59$ in the displayed order.
Within the pair-root family order $(a,b)$ lexicographically and then
$k$ increasingly; within the last family order $(a,b,c)$
lexicographically. The exact certificate partitions the lines into
the following fifteen blocks, labelled by pairs of points:
\begin{center}\small
\begin{tabular}{cl@{\qquad}cl}
\toprule
Pair&Block&Pair&Block\\
\midrule
$12$&$0,1,2,3$&$34$&$4,6,24,26$\\
$56$&$5,7,25,27$&$35$&$8,10,20,22$\\
$46$&$9,11,21,23$&$45$&$12,14,16,18$\\
$36$&$13,15,17,19$&$16$&$28,33,45,48$\\
$26$&$29,32,44,49$&$25$&$30,35,47,50$\\
$15$&$31,34,46,51$&$24$&$36,41,53,56$\\
$14$&$37,40,52,57$&$23$&$38,43,55,58$\\
$13$&$39,42,54,59$&&\\
\bottomrule
\end{tabular}
\end{center}
For every $s\in C$, exact multiplication verifies that its root-line
action preserves these blocks. If $f(s)$ denotes the pair labelling
the block containing its own root, the induced action is
\begin{equation}\label{eq:block}
              \{a,b\}\longmapsto\{f(s)(a),f(s)(b)\},
\end{equation}
where $f(s)$ on the right is a transposition. In particular the five
matrix generators act on the blocks as the five listed
transpositions act on two-element subsets.

This proves a group homomorphism, rather than only a conjugation-table
identity. The exact matrices act on root lines, and the invariant
partition gives an action of $W$ on its fifteen blocks. The action
of $S_6$ on its two-element subsets is faithful. Its image is generated
by the same five block permutations, so identifying these two
permutation images defines $\pi$. The edge graph of the listed
transpositions is connected on six vertices, proving surjectivity.
Equation \eqref{eq:block} identifies the image of every $s\in C$.
Each pair labels one four-element block, giving the four-to-one count.

The exact replay enumerates $46080$ distinct matrices generated by
the five reflections and checks closure under multiplication by them.
It independently enumerates the root-line action, of order $11520$.
Checking rank-one displacement throughout the full matrix group gives
exactly the supplied sixty reflections. Conjugating one reflection
by the five generators reaches all sixty. This proves the reflection
and conjugacy assertions. The first isomorphism theorem now gives
$|\ker\pi|=46080/720=64$.
\end{proof}

\begin{corollary}\label{cor:classrank}
The minimum number of members of $C$ generating $W$ is five.
\end{corollary}
\begin{proof}
The displayed five matrices generate $W$. Any generating set from
$C$ maps to transpositions generating $S_6$. Their edge graph must
be connected on six vertices, requiring at least five edges.
\end{proof}

\section*{Consequences for knot quotients}
A meridional reflection-class quotient means an epimorphism
$\rho:\pi_1(S^3\setminus K)\to W$ taking a meridian to $C$.
All meridians then map into $C$; orientation reversal causes no
change because the reflections are involutions.

\begin{corollary}\label{cor:dominance}
Every meridional reflection-class quotient onto $W$ composes with
$\pi$ to a meridional epimorphism onto the transposition class of
$S_6$. It already gives $\mr(K)\ge5$ through the rank-five Coxeter
group of type $A_5$. It also gives an irreducible five-dimensional
complex KCH representation at meridian parameter $-1$, hence
$\operatorname{ar}(K,\C)\ge5$ for the ordinary $U=1$ augmentation rank.
\end{corollary}
\begin{proof}
Composition preserves surjectivity and takes meridians to
transpositions. Images of any meridional generating set must have a
connected edge graph on six vertices, giving the lower bound five.

The standard representation of $S_6$ on
\[
       U=\{(z_1,\ldots,z_6)\in\C^6:\textstyle\sum_i z_i=0\}
\]
has dimension five, and a transposition has eigenvalues
$-1,1,1,1,1$. It is irreducible: a nonzero invariant subspace
contains a vector with $z_i\ne z_j$, and subtracting its $(ij)$-translate
gives a nonzero multiple of $e_i-e_j$. Its conjugates span $U$.
Composing with the knot-group epimorphism preserves irreducibility,
so this is the asserted KCH representation.
The augmentation conclusion uses the augmentation--KCH correspondence
and rank definition in \cite[Theorem~1.2 and Definition~1.8]{cor17}.
\end{proof}

The natural faithful representation of $G_{31}$ has dimension four,
while its reflection class needs five generators. For knots detected
by this full quotient, that distinction does not give a gap:
the same knot has the preceding five-dimensional KCH representation.
The conclusion concerns this full group and its complete reflection
class. It makes no claim about all higher-rank complex reflection
groups, arbitrary subgroups, modular representations, or completeness
of finite-quotient methods.

\section*{Exact finite support}
The input retains sixty Gaussian-integer root vectors, their rational
Gaussian reflection matrices, five generator indices, the fifteen
blocks and their pair labels. Matrices are encoded by twice their
entries as Gaussian-integer coordinates. The default standard-library
checker reconstructs the reflections from their roots, verifies all
3600 root-line conjugation actions and all $60\cdot15$ block actions,
and independently enumerates the full matrix group, its reflection
set, the root-line group and the target symmetric group.
No Sage, GAP, numerical approximation or earlier computation receipt
is required.

The receipt pins the input, checker, written proof and retained
primary PDFs. The source import names this exact model $G_{31}$;
the finite computations establish the asserted properties of the
explicit model. The homomorphism argument, class-rank implication and
standard-representation proof are given above. This is exact internal
replay, not formal verification or external review. Exploratory braid
and knot examples are outside the certificate.

\begin{thebibliography}{9}
\bibitem{km25}
Y. Konishi and S. Minabe,
\emph{A reduction theorem for good basic invariants of finite complex
reflection groups},
Journal of Algebra \textbf{677} (2025), 327--359,
\url{https://doi.org/10.1016/j.jalgebra.2025.03.048}.
Retained accepted manuscript:
\href{https://arxiv.org/abs/2409.00380v2}{arXiv:2409.00380v2},
Section~6.1, equation~(6.1), printed page~25.
\bibitem{cor17}
C. R. Cornwell, \emph{KCH representations, augmentations, and
$A$-polynomials},
Journal of Symplectic Geometry \textbf{15} (2017), 983--1017.
Theorem numbering follows retained
\href{https://arxiv.org/abs/1310.7526v3}{arXiv:1310.7526v3}.
\end{thebibliography}
\end{document}
